\documentclass[12pt,a4paper]{article}
\usepackage[spanish]{babel}
\usepackage[utf8]{inputenc}
\usepackage[T1]{fontenc}
\usepackage{lmodern}
\usepackage[table]{xcolor}
\usepackage{geometry}
\usepackage{setspace}
\usepackage{longtable}
\usepackage{booktabs}
\usepackage{array}
\usepackage{tabularx}
\usepackage{hyperref}
\usepackage{enumitem}

\geometry{margin=2.4cm}
\setstretch{1.12}

\definecolor{PrimaryBlue}{HTML}{16324F}
\definecolor{SoftGray}{HTML}{F3F6F8}

\title{\Huge\textbf{Sprint Backlog}\\[0.3em]\Large Plataforma Web de Mascotas 3D}
\author{Bryan Quispe}
\date{20 de abril de 2026}

\begin{document}

\begin{titlepage}
\centering
\vspace*{1.8cm}
{\Huge\bfseries\color{PrimaryBlue} Sprint Backlog\\[0.35em]}
{\LARGE Tareas del Sprint con estimación\\[1.2cm]}

\begin{tabular}{p{5cm} p{8cm}}
\rowcolor{SoftGray}
\textbf{Versión} & 1.0 \\
\textbf{Fecha} & 20 de abril de 2026 \\
\textbf{Autor} & Bryan Quispe \\
\textbf{Periodo} & 20 de abril de 2026 al 7 de agosto de 2026 \\
\end{tabular}

\vfill
{\large Documento Scrum para la ejecución del Sprint}
\end{titlepage}

\section{Sprint activo}
\textbf{Sprint 1} - Autenticación y registro básico de mascotas.

\section{Tareas priorizadas}
\begin{longtable}{>{\raggedright\arraybackslash}p{1.8cm} >{\raggedright\arraybackslash}p{6.2cm} >{\raggedright\arraybackslash}p{1.8cm} >{\raggedright\arraybackslash}p{2.2cm} >{\raggedright\arraybackslash}p{1.8cm}}
\toprule
\textbf{ID} & \textbf{Tarea} & \textbf{Historia} & \textbf{Estimación} & \textbf{Estado} \\
\midrule
SB-01 & Crear formulario de registro y validaciones & PB-01 & 5 pts & Pendiente \\
SB-02 & Implementar login con JWT & PB-02 & 5 pts & Pendiente \\
SB-03 & Crear endpoint de alta de mascota & PB-03 & 5 pts & Pendiente \\
SB-04 & Implementar listado de mascotas del usuario & PB-04 & 3 pts & Pendiente \\
SB-05 & Manejo de errores y mensajes consistentes & PB-09 & 3 pts & Pendiente \\
SB-06 & Configurar expiración de sesión por inactividad & PB-06 & 3 pts & Pendiente \\
\bottomrule
\end{longtable}

\section{Capacidad estimada}
\begin{itemize}[leftmargin=*]
    \item Capacidad del Sprint: 24 puntos.
    \item Reserva técnica: 3 puntos.
    \item Capacidad comprometida: 21 puntos.
\end{itemize}

\section{Criterio de control}
\begin{itemize}[leftmargin=*]
    \item Cada tarea debe estar vinculada a una historia priorizada.
    \item No se agregan tareas no refinadas durante el Sprint sin replanificación.
\end{itemize}

\end{document}
